\documentclass[12pt,a4paper]{article}
\usepackage[T1]{fontenc}
\usepackage{amsmath,amssymb,mathrsfs,tikz,times,pifont}
\usepackage{enumitem}
\usepackage{multicol}
\usepackage{lmodern}
\usetikzlibrary{trees}
\newcommand\circitem[1]{%
\tikz[baseline=(char.base)]{
\node[circle,draw=gray, fill=red!55,
minimum size=1.2em,inner sep=0] (char) {#1};}}
\newcommand\boxitem[1]{%
\tikz[baseline=(char.base)]{
\node[fill=cyan,
minimum size=1.2em,inner sep=0] (char) {#1};}}
\setlist[enumerate,1]{label=\protect\circitem{\arabic*}}
\setlist[enumerate,2]{label=\protect\boxitem{\alph*}}
%%%::::::by chnini ameur :::::::%%%
\everymath{\displaystyle}
\usepackage[left=1cm,right=1cm,top=1cm,bottom=1.7cm]{geometry}
\usepackage[colorlinks=true, linkcolor=blue, urlcolor=blue, citecolor=blue]{hyperref}
\usepackage{array,multirow}
\usepackage[most]{tcolorbox}
\usepackage{varwidth}
\usepackage{float} %pour utiliser l'option [H] qui force l'image à apparaître exactement à l'endroit où elle est placée dans le code.
\tcbuselibrary{skins,hooks}
\usetikzlibrary{patterns}
%%%::::::by chnini ameur :::::::%%%
\newtcolorbox{exa}[2][]{enhanced,breakable,before skip=2mm,after skip=5mm,
colback=yellow!20!white,colframe=black!20!blue,boxrule=0.5mm,
attach boxed title to top left ={xshift=0.6cm,yshift*=1mm-\tcboxedtitleheight},
fonttitle=\bfseries,
title={#2},#1,
% varwidth boxed title*=-3cm,
boxed title style={frame code={
\path[fill=tcbcolback!30!black]
([yshift=-1mm,xshift=-1mm]frame.north west)
arc[start angle=0,end angle=180,radius=1mm]
([yshift=-1mm,xshift=1mm]frame.north east)
arc[start angle=180,end angle=0,radius=1mm];
\path[left color=tcbcolback!60!black,right color = tcbcolback!60!black,
middle color = tcbcolback!80!black]
([xshift=-2mm]frame.north west) -- ([xshift=2mm]frame.north east)
[rounded corners=1mm]-- ([xshift=1mm,yshift=-1mm]frame.north east)
-- (frame.south east) -- (frame.south west)
-- ([xshift=-1mm,yshift=-1mm]frame.north west)
[sharp corners]-- cycle;
},interior engine=empty,
},interior style={top color=yellow!5}}
%%%%%%%%%%%%%%%%%%%%%%%
\usepackage{fancyhdr}
\usepackage{eso-pic}         % Pour ajouter des éléments en arrière-plan
% Commande pour ajouter du texte en arrière-plan
\usepackage{tkz-tab}
\usepackage{tikz}
\usepackage[siunitx, european, RPvoltages]{circuitikz}
\AddToShipoutPicture{
    \AtTextCenter{%
        \makebox[0pt]{\rotatebox{80}{\textcolor[gray]{0.7}{\fontsize{5cm}{5cm}\selectfont PGB}}}
    }
}
\usepackage{lastpage}
\fancyhf{}
\pagestyle{fancy}
\renewcommand{\footrulewidth}{1pt}
\renewcommand{\headrulewidth}{0pt}
\renewcommand{\footruleskip}{10pt}
\fancyfoot[R]{
\color{blue}\ding{45}\ \textbf{2026}
}
\fancyfoot[L]{
\color{blue}\ding{45}\ \textbf{Prof:M. BA}
}
\cfoot{\bf
\thepage /
\pageref{LastPage}}
% Création du compteur pour les exercices
\newcounter{exercice}
\renewcommand{\theexercice}{\arabic{exercice}}  % Définit l'affichage du compteur en chiffres arabes

% Définir la commande \exo
\newcommand{\exo}{\refstepcounter{exercice}\textbf{Exercice \theexercice} }
\begin{document}
\renewcommand{\arraystretch}{1.5}
\renewcommand{\arrayrulewidth}{1.2pt}
\begin{tikzpicture}[overlay,remember picture]
    \node[draw=blue,line width=1.2pt,fill=purple,text=blue,inner sep=3mm,rounded corners,pattern=dots]at ([yshift=-2.5cm]current page.north) {\begingroup\setlength{\fboxsep}{0pt}\colorbox{white}{\begin{tabular}{|*1{>{\centering \arraybackslash}p{0.28\textwidth}} |*2{>{\centering \arraybackslash}p{0.2\textwidth}|} *1{>{\centering \arraybackslash}p{0.19\textwidth}|} }
                \hline
                \multicolumn{3}{|c|}{$\diamond$$\diamond$$\diamond$\ \textbf{Lycée de Dindéfélo}\ $\diamond$$\diamond$$\diamond$ } & \textbf{A.S. : 2026/2027}                                              \\ \hline
                \textbf{Matière: Mathématiques}                                                                                    & \textbf{Niveau : 1$^{er}$}\textbf{S} & \textbf{Date: 18/10/2026} & \textbf{} \\ \hline
                \multicolumn{4}{|c|}{\parbox[c]{10cm}{\begin{center}
                                                                  \textbf{{\Large\sffamily Equations-Inéquations Irrationnelles et Système}}
                                                              \end{center}}}                                                                                                        \\ \hline
            \end{tabular}}\endgroup};
\end{tikzpicture}
\vspace{4cm}

\fbox{\textbf{\exo}}

\begin{enumerate}
\item Résoudre dans \( \mathbb{R} \) les équations et inéquations suivantes :

\[ \text{a)} \sqrt{2x-1}-\sqrt{x-5} = 0    \quad\quad;\quad\quad   \text{b)} \sqrt{4x^{2}+x+5}=2x+1 \]

\[ \text{c)} \sqrt{-x^{2}+x+1} \leq x-2    \quad\quad;\quad\quad   \text{d)} \sqrt{x^{2}+6x+6} \geq 2x+1 \]

\[ \text{e)} \sqrt{2x^{2}-3x-2} = 2x-4    \quad\quad;\quad\quad   \text{f)} \sqrt{4-x} = \sqrt{x^{2}+2x-6} \]

\[ \text{g)} \sqrt{x^{2}+4x} < \sqrt{x^{2}-5x+4}    \quad\quad;\quad\quad   \text{h)} \sqrt{x^{2}-x-2} > x-1 \]

\end{enumerate}

\fbox{\textbf{\exo}}

\begin{enumerate}
\item Soit \( m \) un nombre réel. On considère l’équation d’inconnue \( x \) :

\( (E) \) \( (m+3)x^{2} +2mx + m + 5 = 0\)
\item[a)] Discuter suivant les valeurs de \( m \), l’existence, le nombre et le signe des racines de \( (E) \) 

Dans le cas où \( (E) \) admet deux racines distinctes \( x_{2} \) et \( x_{2} \) , peut-on déterminer \( m \) pour que:

\( ( 2x_{1} - 1 ) (2x_{2}-1) = 6 \) puis \( x_{1}^{2} + x_{2}^{2} = 2 \)

\item[b)]Trouver une relation indépendante de \( m \) entre\( x_{2} \) et \( x_{2} \)
\item[c)]Utiliser cette relation pour trouver les racines doubles de \( (E) \), lorsqu’elles existent
\end{enumerate}

\fbox{\textbf{\exo}}

\begin{enumerate}[label=\Roman*.]
    \item Soit l'équation $(E) : x^2 - (2m + 1)x + m^2 - 2 = 0$
    \begin{enumerate}[label=\arabic*.]
        \item Discuter valeurs de $m$, les solutions de l'équation $(E)$
        \item Dans le cas où $(E)$ admet deux solutions $x'$ et $x''$, déterminer $m$ tel que :
        \begin{enumerate}[label=\alph*.]
            \item $x' + x'' = 35$
            \item $|x' - x''| = 5$
        \end{enumerate}
    \end{enumerate}
    
    \item On considère l'équation $(E_m) : 3(-m + 2)x^2 + (3m + 9)x + 9m + 3 = 0$.
    \begin{enumerate}[label=\arabic*.]
        \item Déterminer $m$ pour que $-1$ soit racine de $(E_m)$.
        \item Résoudre $(E_m)$ pour $m = 2$.
        \item On suppose que $m \neq 2$.
        \begin{enumerate}[label=\alph*.]
            \item Montrer que le discriminant $\Delta = 9(13m^2 - 14m + 1)$
            \item Pour quelles valeurs de $m$ l'équation n'admet pas de solutions.
            \item Pour quelles valeurs de $m$ l'équation admet des solutions distinguées.
            \item Dans le cas où les racines existent.
            \begin{enumerate}[label=\roman*.]
                \item Trouver une relation indépendante de $m$ entre les racines.
                \item Trouver alors les racines doubles de $(E_m)$.
            \end{enumerate}
        \end{enumerate}
    \end{enumerate}
\end{enumerate}

\fbox{\textbf{\exo}}

\begin{enumerate}
\item Résoudre dans \( \mathbb{R}^{3} \) par le pivot de Gauss le système \( \begin{cases} 27x + 9y + 3z = -18\\ -x + y -z = 2\\ x + y + z = -8 \end{cases}\)

\item Résoudre dans \( \mathbb{R}^{3} \) par le pivot de Gauss le système \( \begin{cases} 4x - 2y + z = 6\\ 8x + 6y -3z = 22\\ 2x + 4y + 7z = 44 \end{cases}\)
\item En déduire la résolution dans \( \mathbb{R}^{3} \) du système \( \begin{cases} \frac{4}{x+3} - \frac{2}{y-2} + \frac{1}{z} = 6\\ \frac{8}{x+3} + \frac{6}{y-2} -\frac{3}{z} = 22\\ \frac{2}{x+3} + \frac{4}{y-2} + \frac{7}{z} = 44 \end{cases}\)
\end{enumerate}

\fbox{\textbf{\exo}}

\begin{enumerate}
\item Soit le polynôme de degré 3 tel que \( P(0) = 8 \) ; \( P(1) = 0 \) ; \( P(3) = -10 \) et \( P(-1) = 10 \).
\begin{enumerate}
\item Montrer que \( P(x) = x^{3} - 3x^{2} - 6x + 8 \).
\item Résoudre dans \( \mathbb{R} \) \( P(x) = 0 \).
\item Résoudre dans \( \mathbb{R} \) \( P(\sqrt{x+3}-x) = 0\).
\end{enumerate}
\end{enumerate}

\fbox{\textbf{\exo}}

\begin{enumerate}
    \item Résoudre dans $ \mathbb{R} $, les équations suivantes :
    \begin{enumerate}
        \item $\sqrt{x^2 - 5x} = x$
        \item $\sqrt{x^2 - 5x + 6} \geq x - 2$
        \item $\sqrt{2x^2 - x} < 2x - 3$
    \end{enumerate}
    
    \item 
    \begin{enumerate}
        \item[a)] Résoudre dans $\mathbb{R}$, l'équation suivante : $2x^2 - 4x - 30 = 0$
        \item[b)] En déduire l’ensemble des solutions de l’équation : $2(x^2 - 1)^2 - 4|x^2 - 1| - 30 = 0$
    \end{enumerate}
    
    \item Résoudre le système suivant par la méthode du pivot de Gauss :
    \[
    \begin{cases}
        5x - 4y - z = 3 \\
        -x + 2y - 3z = 7 \\
        x - 3y + 2z = 5
    \end{cases}
    \]
\end{enumerate}

\fbox{\textbf{\exo}}

Soit l'équation $(E_m) : (m - 2)x^2 + (2m + 2)x + 10m - 14 = 0 ; m \in \mathbb{R}$

\begin{enumerate}
    \item Discuter suivant les valeurs de $m$ le nombre de racines de $(E_m)$
    \item Pour quelle(s) valeur(s) de $m$, les solutions $x_1$ et $x_2$ de $(E_m)$ :
    \begin{enumerate}
        \item sont de signes contraires ?
        \item sont strictement positives ?
    \end{enumerate}
\end{enumerate}

\fbox{\textbf{\exo}}

\begin{enumerate}
    \item Résoudre dans $\mathbb{R}$ les équations suivantes.
    \begin{enumerate}
        \item $\sqrt{-x^2 + 5x + 9} = \sqrt{x - 3}$
        \item $\sqrt{x^2 - x - 2} = x + 1$
    \end{enumerate}
    
    \item Résoudre dans $\mathbb{R}$ les inéquations suivantes.
    \begin{enumerate}
        \item $\sqrt{(x + 3)(x - 1)} \leq x - 3$
        \item $\sqrt{-5x + 1} \geq \sqrt{x^2 - 3x + 2}$
        \item $x^4 + x^2 - 12 \leq 0$
    \end{enumerate}
    
    \item Résoudre dans $\mathbb{R}^2$ le système suivant
    \[
    \begin{cases}
        x^3 + y^3 = 9 \\
        \frac{x}{y} + \frac{y}{x} = \frac{5}{2}
    \end{cases}
    \]
    
    \item Résoudre dans $\mathbb{R}^3$ par la méthode du Pivot de Gauss le système
    \[
    \begin{cases}
        x - 2y + z = 6 \\
        -2x + y - z = -6 \\
        3x - y - 2z = -2
    \end{cases}
    \]
\end{enumerate}

\fbox{\textbf{\exo}}

On donne \( P(x) = 2x^3 + ax^2 + bx - 6 \) où \( a \) et \( b \) sont des réels.
\begin{enumerate}
    \item Trouver les réels \( a \) et \( b \) pour que le polynôme \( P(x) \) soit divisible par \( x^2 - x - 2 \).
    \item En déduire une factorisation de \( P(x) \).
    \item Résoudre dans \( \mathbb{R} \) l'équation \( P(x) = 0 \).
    \item En déduire les solutions de l'équation \( P(x^2 - 2) = 0 \).
    \item Résoudre dans \( \mathbb{R} \) l'inéquation \( P(x) \leq 0 \).
    \item En déduire les solutions de l'inéquation \( P(3 - 2x) \leq 0 \).
    \item Résoudre dans \( \mathbb{R} \) l'inéquation \( \frac{2x^3 + x^2 - 7x - 6}{-x + 1} > 0 \).
\end{enumerate}

\end{document}